\documentclass[11pt]{article}

\usepackage{amsmath,amssymb}
\usepackage{xurl}
\usepackage[hidelinks]{hyperref}
\setlength{\emergencystretch}{2em}

% Shared commands for notes in docs/paper-gaps/.
%
% Two halves.  The link commands below are project identity: OWNER/REPO is the
% placeholder that scripts/bootstrap_project.py replaces with this project's
% GitHub slug and Pages site.  Everything after them is NOTATION, and notation
% belongs to the source paper: the definitions here are an example set, and a
% project replaces them with the macros of its own paper's style file.  The one
% rule that never changes: a command defined here must mean exactly what the
% source paper's macro means, or a note silently says something else.

\newcommand{\Lean}{\textsc{Lean}}
\newcommand{\leanid}[1]{\textsf{#1}}
\newcommand{\ghissue}[1]{\href{https://github.com/OWNER/REPO/issues/#1}{GitHub issue \##1}}
\newcommand{\ghpr}[1]{\href{https://github.com/OWNER/REPO/pull/#1}{PR \##1}}
% Local issue tracker (issues/NNNN-*.md in this repository); no remote URL.
\newcommand{\localissue}[1]{issue \##1 (\path{issues/})}
\newcommand{\doclink}[1]{\href{https://OWNER.github.io/REPO/docs/#1.html}{\path{#1}}}

% --- Notation: replace with the source paper's own macros --------------------
% \F, \N, \C, \R, \tr, \ot, \eps, \E, \bx, \bu, \bv, \by

\newcommand{\F}{\mathbb{F}}
\newcommand{\N}{\mathbb{N}}
\newcommand{\C}{\mathbb{C}}
\newcommand{\R}{\mathbb{R}}
\newcommand{\tr}{\mathrm{tr}}
\newcommand{\eps}{\epsilon}
\newcommand{\ot}{\otimes}
\newcommand{\E}{\mathop{\bf E\/}}

% bold random variables
\newcommand{\bu}{\boldsymbol{u}}
\newcommand{\bv}{\boldsymbol{v}}
\newcommand{\bx}{{\boldsymbol{x}}}
\newcommand{\by}{\boldsymbol{y}}

% T1 so literal < and > in \texttt placeholders typeset as themselves
\usepackage[T1]{fontenc}

% bra-ket
\usepackage{braket}

% --- Example structured spaces ---
\newcommand{\polyfunc}[3]{\mathcal{P}(#1, #2, #3)}
\newcommand{\polymeas}[3]{\mathrm{PolyMeas}(#1, #2, #3)}
\newcommand{\polysub}[3]{\mathrm{PolySub}(#1, #2, #3)}

% Approximation relations (paper uses \approx_\eps and \simeq_\zeta)
\newcommand{\approxeps}{\approx_{\varepsilon}}
\newcommand{\simgeq}{\simeq_{\zeta}}


\renewcommand{\ghissue}[1]{%
  \href{https://github.com/Dengnifer/MIPStarRE-A/issues/#1}{GitHub issue \##1}}

\title{Raw Pauli Effects in the Soundness Conclusion}
\author{}
\date{2026-09-21}

\begin{document}
\maketitle

\section{At a glance}

\begin{itemize}
  \item \textbf{Difficulty.} Medium.  The repair must control a positive
  malformed-answer effect on the extracted ideal state for both players.
  \item \textbf{Estimated weight.} Two rejection-mass estimates, two
  completed-to-raw distance estimates, and one scalar-absorption step.
  \item \textbf{Mathlib/project split.} Matrix positivity, contraction, and
  norm inequalities come from Mathlib and the quantum API; the oriented-edge
  rejection bounds and Pauli-answer parsing are project-local.
  \item \textbf{Key mathematical inputs.} Completion changes only outcome
  zero, malformed Pauli answers are rejected on a fixed Pauli--point edge,
  and the robustness family absorbs linear and squared error terms.  The
  repair is tracked by \ghissue{665} and its documentation by \ghissue{666}.
\end{itemize}

\section{Key theorem forms}

Theorem \path{thm:pauli} compares, for each $W\in\{X,Z\}$ and
$u\in\F_q^M$, the effect $A^{(\mathrm{Pauli},W)}_u$ or
$B^{(\mathrm{Pauli},W)}_u$ directly with the corresponding ideal Pauli
projector.  The same raw prescribed-answer effects occur in
\path{cor:pauli-binary} of \cite{Ji2020MIPStar}.%
\footnote{See
\path{references/qpbt-paper/08_classical_and_quantum_low_degree_tests.tex},
lines~1426--1487, and the closing proof in
\path{references/qpbt-paper/14_analysis_of_the_pauli_basis_test.tex},
lines~1862--1876.}

\section{Formal representation discrepancy}

The formal game has one sum-type answer alphabet for all question types.
Earlier soundness declarations first postprocessed a Pauli-question
measurement by sending every malformed answer to the register value zero.
For $u\ne0$ this completed effect equals the raw effect attached to
\texttt{pauliOutcome $u$}.  At zero it is instead
\[
  M^{\mathrm{completed}}_0
  = M_{\texttt{pauliOutcome }0}+M_{\mathrm{malformed}}.
\]
Consequently the old headline distance compared the completed family, not the
raw family in the paper.  Rejection of malformed answers did not make
$M_{\mathrm{malformed}}$ identically zero, so this was not a definitional
identification.

\section{Repair}

For Alice, rejection on the oriented Pauli-to-point edge bounds the state
mass of $M_{\mathrm{malformed}}$ by
$|\mathrm{PauliEdge}|\,\varepsilon$.  The reverse point-to-Pauli edge gives
the same bound for Bob.  If the transported strategy state is within
$\delta$ of the ideal state, contractivity gives
\[
  \|\widetilde M_{\mathrm{malformed}}\ket{\mathrm{ideal}}\|^2
  \leq 2\delta^2+2|\mathrm{PauliEdge}|\,\varepsilon.
\]
Since completion changes only outcome zero, the squared triangle inequality
then yields, for each player,
\[
  d_{\mathrm{raw}}
  \leq 2d_{\mathrm{completed}}+4\delta^2
      +4|\mathrm{PauliEdge}|\,\varepsilon.
\]
On $0\leq\varepsilon\leq1$, the soundness function dominates
$\varepsilon$; the existing scalar absorption estimates therefore absorb all
three terms by enlarging only the existential prefactor $a$.  The usual
monotonicity argument covers $\varepsilon>1$.

The declarations
\leanid{raw\_pauli\_operator\_distanceA\_le\_completed} and
\leanid{raw\_pauli\_operator\_distanceB\_le\_completed} implement the two
transfer estimates, and
\leanid{exists\_arbitrary\_strategy\_raw\_isometry\_bounds} performs the
scalar absorption.  The public theorems \leanid{pauli\_soundness} and
\leanid{pauli\_soundness\_qubit} now use the raw effects directly.  The
completed-family distances remain as internal support for extraction and
Naimark dilation.

\section{Conclusion}

The paper statement is unchanged and requires no additional hypothesis.  The
formalization now proves its raw-effect conclusion for both players with the
printed form of $\delta_{\mathrm{qld}}$.  The discrepancy and repair are
tracked by \ghissue{665}; documentation follow-up is tracked by
\ghissue{666}.

\bibliographystyle{alpha}
\bibliography{references}

\end{document}
